\documentclass[11pt]{article}

\usepackage[margin=1in]{geometry}
\usepackage{amsmath,amssymb,amsthm,mathtools}
\usepackage{hyperref}
\usepackage{enumitem}

\hypersetup{colorlinks=true,linkcolor=blue,citecolor=blue,urlcolor=blue}

\newtheorem{definition}{Definition}
\newtheorem{assumption}{Assumption}
\newtheorem{theorem}{Theorem}
\newtheorem{lemma}{Lemma}
\newtheorem{corollary}{Corollary}
\newtheorem{remark}{Remark}

\newcommand{\R}{\mathbb{R}}
\newcommand{\ip}[2]{\left\langle #1,#2\right\rangle}
\newcommand{\ipx}[2]{\left\langle #1,#2\right\rangle_x}
\newcommand{\norm}[1]{\left\|#1\right\|}
\newcommand{\dd}{\,d}

\title{Level-Set Estimates from Hessian Stability and the Asymmetry Residual}
\author{}
\date{}

\begin{document}
\maketitle

\section{Setup and notation}

Let \(\phi\) be a strictly convex \(C^3\) function on a convex domain
\(\Omega\subseteq\R^n\). Write
\[
    g(x):=\nabla \phi(x),
    \qquad
    H(x):=\nabla^2\phi(x).
\]
The Hessian metric defines the local primal and dual norms
\[
    \|u\|_x^2:=u^\top H(x)u,
    \qquad
    \|w\|_{x,*}^2:=w^\top H(x)^{-1}w.
\]
Let \(d_R\) denote the Riemannian distance induced by the Hessian metric. Define
\[
    V(x,y):=\ip{x-y}{g(x)-g(y)},
\]
\[
    d(x,y):=
    \|x-y\|_x^2+
    \|g(x)-g(y)\|_{x,*}^2,
\]
and
\[
    U(x,y):=(x-y)-H(x)^{-1}\bigl(g(x)-g(y)\bigr).
\]
We use \(\ipx{a}{b}:=a^\top H(x)b\). Thus \(\|a\|_x^2=\ipx{a}{a}\).

\begin{definition}[Hessian stability]
Let \(\rho:[0,\infty)\to[1,\infty)\) be nondecreasing with \(\rho(0)=1\). We say that the Hessian metric is \(\rho\)-stable if, for all \(p,q\in\Omega\),
\[
    \rho(d_R(p,q))^{-1}H(p)
    \preceq
    H(q)
    \preceq
    \rho(d_R(p,q))H(p).
\]
We write
\[
    K_\rho(r):=\sqrt{\rho(r)}.
\]
The exponential case is \(\rho(r)=e^{Lr}\).
\end{definition}

\begin{lemma}[Universal distance-divergence lower bound]\label{lem:distance-V}
For all \(x,y\in\Omega\),
\[
    d_R(x,y)^2\le V(x,y).
\]
\end{lemma}

\begin{proof}
Let \(h:=y-x\), and consider the chord \(c(t):=x+th\), \(t\in[0,1]\). Since \(\Omega\) is convex, the chord lies in \(\Omega\). Its Riemannian length is
\[
    \ell(c)=\int_0^1 \sqrt{h^\top H(c(t))h}\,dt.
\]
By Cauchy--Schwarz,
\[
    \ell(c)^2
    \le
    \int_0^1 h^\top H(c(t))h\,dt.
\]
The right-hand side equals
\[
    h^\top\left(\int_0^1 H(x+th)\,dt\right)h
    =
    h^\top(g(y)-g(x))
    =
    V(x,y).
\]
Since \(d_R(x,y)\le \ell(c)\), the claim follows.
\end{proof}

\section{Two estimates from Hessian stability}

The next two lemmas are independent consequences of the same stability assumption. The first controls the asymmetry residual \(U\). The second controls \(V\) by geodesic distance.

\begin{lemma}[Asymmetry residual bound]\label{lem:U-bound}
Assume \(\rho\)-Hessian stability. Define
\[
    J_\rho(R):=
    \int_0^R
    \left(
        K_\rho(s)-K_\rho(s)^{-1}
    \right)ds.
\]
Then
\[
    \|U(x,y)\|_x
    \le
    J_\rho(d_R(x,y)).
\]
Equivalently,
\[
    \|U(x,y)\|_x
    \le
    \int_0^{d_R(x,y)}
    \left(
        \sqrt{\rho(s)}
        -
        \frac{1}{\sqrt{\rho(s)}}
    \right)ds.
\]
\end{lemma}

\begin{proof}
Let \(R:=d_R(x,y)\), and let \(\gamma:[0,R]\to\Omega\) be a unit-speed geodesic from \(x\) to \(y\):
\[
    \gamma(0)=x,
    \qquad
    \gamma(R)=y,
    \qquad
    \|\gamma'(s)\|_{\gamma(s)}=1.
\]
By the fundamental theorem of calculus,
\[
    y-x=\int_0^R \gamma'(s)\,ds,
\]
and
\[
    g(y)-g(x)
    =
    \int_0^R H(\gamma(s))\gamma'(s)\,ds.
\]
Therefore
\[
\begin{aligned}
    U(x,y)
    &=
    (x-y)-H(x)^{-1}(g(x)-g(y)) \\
    &=
    H(x)^{-1}(g(y)-g(x))-(y-x) \\
    &=
    \int_0^R
    \left[
        H(x)^{-1}H(\gamma(s))-I
    \right]\gamma'(s)\,ds.
\end{aligned}
\]

We bound the integrand in the \(x\)-norm. Define
\[
    M_s:=H(x)^{-1/2}H(\gamma(s))H(x)^{-1/2},
    \qquad
    u_s:=H(x)^{1/2}\gamma'(s).
\]
Then
\[
    \left\|
    \left[H(x)^{-1}H(\gamma(s))-I\right]\gamma'(s)
    \right\|_x
    =
    \|(M_s-I)u_s\|_2.
\]
By stability and \(d_R(x,\gamma(s))\le s\),
\[
    \rho(s)^{-1}I\preceq M_s\preceq \rho(s)I.
\]
Also, since \(\gamma\) is unit speed,
\[
    u_s^\top M_su_s
    =
    \gamma'(s)^\top H(\gamma(s))\gamma'(s)
    =1.
\]
Diagonalize \(M_s\). Its eigenvalues \(\lambda_i\) lie in \([\rho(s)^{-1},\rho(s)]\). Writing \(u_s=\sum_i u_i e_i\), we have
\[
    \sum_i\lambda_i u_i^2=1,
\]
and
\[
    \|(M_s-I)u_s\|_2^2
    =
    \sum_i(\lambda_i-1)^2u_i^2.
\]
For every \(\lambda\in[\rho(s)^{-1},\rho(s)]\),
\[
    \frac{|\lambda-1|}{\sqrt\lambda}
    \le
    \sqrt{\rho(s)}-\frac{1}{\sqrt{\rho(s)}}.
\]
Thus
\[
    (\lambda_i-1)^2
    \le
    \left(
        \sqrt{\rho(s)}-\frac{1}{\sqrt{\rho(s)}}
    \right)^2\lambda_i.
\]
It follows that
\[
\begin{aligned}
    \|(M_s-I)u_s\|_2^2
    &\le
    \left(
        \sqrt{\rho(s)}-\frac{1}{\sqrt{\rho(s)}}
    \right)^2
    \sum_i\lambda_i u_i^2 \\
    &=
    \left(K_\rho(s)-K_\rho(s)^{-1}\right)^2.
\end{aligned}
\]
Hence
\[
    \left\|
    \left[H(x)^{-1}H(\gamma(s))-I\right]\gamma'(s)
    \right\|_x
    \le
    K_\rho(s)-K_\rho(s)^{-1}.
\]
Integrating gives
\[
\begin{aligned}
    \|U(x,y)\|_x
    &\le
    \int_0^R
    \left\|
    \left[H(x)^{-1}H(\gamma(s))-I\right]\gamma'(s)
    \right\|_x\,ds \\
    &\le
    \int_0^R\left(K_\rho(s)-K_\rho(s)^{-1}\right)ds
    =J_\rho(R).
\end{aligned}
\]
\end{proof}

\begin{lemma}[Jeffreys divergence bound]\label{lem:V-bound}
Assume \(\rho\)-Hessian stability. Define
\[
    \Psi_\rho(R):=
    2\int_0^R (R-r)K_\rho(r)\,dr.
\]
Then
\[
    V(x,y)
    \le
    \Psi_\rho(d_R(x,y)).
\]
Equivalently,
\[
    V(x,y)
    \le
    2\int_0^{d_R(x,y)}
    \left(d_R(x,y)-r\right)\sqrt{\rho(r)}\,dr.
\]
\end{lemma}

\begin{proof}
Let \(R:=d_R(x,y)\), and let \(\gamma:[0,R]\to\Omega\) be a unit-speed geodesic from \(x\) to \(y\). Then
\[
    y-x=\int_0^R \gamma'(a)\,da,
\]
and
\[
    g(y)-g(x)
    =
    \int_0^R H(\gamma(b))\gamma'(b)\,db.
\]
Thus
\[
\begin{aligned}
    V(x,y)
    &=
    \ip{x-y}{g(x)-g(y)} \\
    &=
    \ip{y-x}{g(y)-g(x)} \\
    &=
    \int_0^R\int_0^R
    \left\langle
        \gamma'(a),
        H(\gamma(b))\gamma'(b)
    \right\rangle
    da\,db.
\end{aligned}
\]
For each \(a,b\), Cauchy--Schwarz in the metric \(H(\gamma(b))\) gives
\[
    \left\langle
        \gamma'(a),
        H(\gamma(b))\gamma'(b)
    \right\rangle
    \le
    \|\gamma'(a)\|_{\gamma(b)}
    \|\gamma'(b)\|_{\gamma(b)}.
\]
Since \(\gamma\) is unit speed, \(\|\gamma'(b)\|_{\gamma(b)}=1\). By stability,
\[
    H(\gamma(b))
    \preceq
    \rho(d_R(\gamma(a),\gamma(b)))H(\gamma(a)).
\]
Since \(\gamma\) is unit speed,
\[
    d_R(\gamma(a),\gamma(b))\le |a-b|.
\]
Therefore
\[
    \|\gamma'(a)\|_{\gamma(b)}
    \le
    K_\rho(|a-b|)\|\gamma'(a)\|_{\gamma(a)}
    =K_\rho(|a-b|).
\]
Thus
\[
    V(x,y)
    \le
    \int_0^R\int_0^R K_\rho(|a-b|)\,da\,db.
\]
Using the interval-pair identity
\[
    \int_0^R\int_0^R f(|a-b|)\,da\,db
    =
    2\int_0^R(R-r)f(r)\,dr,
\]
with \(f=K_\rho\), we obtain
\[
    V(x,y)
    \le
    2\int_0^R(R-r)K_\rho(r)\,dr
    =
    \Psi_\rho(R).
\]
\end{proof}


\begin{corollary}[Exponential specialization]\label{cor:exp-stability-formulas}
If \(\rho(r)=e^{Lr}\), then
\[
    J_\rho(R)=\frac{4}{L}\left(\cosh\left(\frac{LR}{2}\right)-1\right)
\]
and
\[
    \Psi_\rho(R)=\frac{8}{L^2}\left(e^{LR/2}-1\right)-\frac{4}{L}R.
\]
For standard self-concordance \((L=2)\),
\[
    \|U(x,y)\|_x\le 2\left(\cosh d_R(x,y)-1\right)
\]
and
\[
    V(x,y)\le 2\left(e^{d_R(x,y)}-1-d_R(x,y)\right).
\]
\end{corollary}

\section{The parallelogram identity and energy control}

The next identity removes the need for a separate direct bound on \(d(x,y)\). It expresses \(d(x,y)\) exactly in terms of \(V(x,y)\) and \(U(x,y)\).

\begin{lemma}[Parallelogram identity]\label{lem:parallelogram}
Let
\[
    p:=x-y,
    \qquad
    q:=H(x)^{-1}(g(x)-g(y)).
\]
Then
\[
    V(x,y)=\ipx{p}{q},
    \qquad
    d(x,y)=\|p\|_x^2+\|q\|_x^2,
    \qquad
    U(x,y)=p-q.
\]
Consequently,
\[
    d(x,y)=2V(x,y)+\|U(x,y)\|_x^2.
\]
\end{lemma}

\begin{proof}
The identities for \(V\), \(d\), and \(U\) follow immediately from the definitions. Then
\[
    \|U(x,y)\|_x^2
    =
    \|p-q\|_x^2
    =
    \|p\|_x^2+\|q\|_x^2-2\ipx{p}{q}
    =
    d(x,y)-2V(x,y).
\]
Rearranging proves the claim.
\end{proof}

\begin{corollary}[Energy control from the asymmetry residual]\label{cor:energy-control}
Suppose an update satisfies
\[
    E\le d(x,y).
\]
If
\[
    \|U(x,y)\|_x\le A(V(x,y))
\]
for some scalar function \(A\), then
\[
    E\le 2V(x,y)+A(V(x,y))^2.
\]
\end{corollary}

\begin{proof}
By Lemma~\ref{lem:parallelogram},
\[
    d(x,y)=2V(x,y)+\|U(x,y)\|_x^2.
\]
Combining this with \(E\le d(x,y)\) and \(\|U(x,y)\|_x\le A(V(x,y))\) gives the result.
\end{proof}


\section{Cubic control and geodesic acceleration}

The level-set estimates below require control of the geodesic acceleration.  For
Hessian metrics this is an infinitesimal consequence of Hessian stability, provided
that the stability modulus has a finite logarithmic slope at the origin.

\begin{definition}[Compatible stability modulus]\label{def:compatible-modulus}
A stability modulus \(\rho\) is called \emph{compatible} if the one-sided limit
\[
    L_\rho
    :=
    \lim_{r\downarrow0}\frac{\log \rho(r)}{r}
\]
exists and is finite.  In estimates below, the limit may be replaced by the
corresponding \(\limsup\); assuming the limit exists simply avoids notational
clutter.
\end{definition}

\begin{lemma}[Cubic bound from compatible Hessian stability]\label{lem:rho-to-cubic}
Assume the Hessian metric is \(\rho\)-stable and that \(\rho\) is compatible with
constant \(L_\rho\).  Then, for all \(x\in\Omega\) and all vectors \(a,b\),
\[
    |D^3\phi(x)[a,b,b]|
    \le
    L_\rho\|a\|_x\|b\|_x^2.
\]
\end{lemma}

\begin{proof}
If \(a=0\) there is nothing to prove.  Otherwise set
\(\hat a:=a/\|a\|_x\).  For sufficiently small \(\varepsilon>0\), the point
\[
    y_\varepsilon:=x+\varepsilon \hat a
\]
lies in \(\Omega\), and the Riemannian distance satisfies
\[
    d_R(x,y_\varepsilon)=\varepsilon+o(\varepsilon).
\]
By compatibility,
\[
    \rho(d_R(x,y_\varepsilon))
    =
    1+L_\rho\varepsilon+o(\varepsilon).
\]
Hessian stability gives
\[
    D^2\phi(y_\varepsilon)[b,b]
    \le
    \rho(d_R(x,y_\varepsilon))D^2\phi(x)[b,b].
\]
Using \(C^3\)-smoothness,
\[
    D^2\phi(y_\varepsilon)[b,b]
    =
    D^2\phi(x)[b,b]
    +
    \varepsilon D^3\phi(x)[\hat a,b,b]
    +o(\varepsilon).
\]
Combining the previous two displays and cancelling \(D^2\phi(x)[b,b]=\|b\|_x^2\)
gives
\[
    D^3\phi(x)[\hat a,b,b]
    \le
    L_\rho\|b\|_x^2.
\]
Repeating the same argument with \(-\hat a\) in place of \(\hat a\) gives
\[
    -D^3\phi(x)[\hat a,b,b]
    \le
    L_\rho\|b\|_x^2.
\]
Thus
\[
    |D^3\phi(x)[\hat a,b,b]|
    \le
    L_\rho\|b\|_x^2.
\]
Multiplying by \(\|a\|_x\) proves the claim.
\end{proof}

\begin{lemma}[Acceleration from a cubic bound]\label{lem:acceleration}
Suppose that, for some \(L\ge0\),
\[
    |D^3\phi(x)[a,b,b]|\le L\|a\|_x\|b\|_x^2
\]
for all \(x,a,b\). Let \(x(t)\) be a Hessian geodesic with constant energy
\[
    E:=\|\dot x(t)\|_{x(t)}^2.
\]
Then
\[
    \|\ddot x(t)\|_{x(t)}\le \frac{L}{2}E.
\]
In particular, standard self-concordance \((L=2)\) gives
\[
    \|\ddot x(t)\|_{x(t)}\le E.
\]
\end{lemma}

\begin{proof}
The Hessian geodesic equation is
\[
    H(x(t))\ddot x(t)
    =-\frac12D^3\phi(x(t))[\dot x(t),\dot x(t),\cdot].
\]
Therefore, for any \(w\),
\[
    \ip{\ddot x(t)}{w}_{x(t)}
    =-\frac12D^3\phi(x(t))[\dot x(t),\dot x(t),w].
\]
By symmetry of the cubic tensor and the assumed cubic bound,
\[
    \left|D^3\phi(x(t))[\dot x(t),\dot x(t),w]\right|
    \le L\|w\|_{x(t)}\|\dot x(t)\|_{x(t)}^2
    =L\|w\|_{x(t)}E.
\]
Taking the supremum over \(\|w\|_{x(t)}\le1\) gives the result.
\end{proof}

\begin{corollary}[Acceleration from compatible Hessian stability]\label{cor:rho-acceleration}
If the Hessian metric is \(\rho\)-stable with compatible modulus \(\rho\), then every
Hessian geodesic satisfies
\[
    \|\ddot x(t)\|_{x(t)}
    \le
    \frac{L_\rho}{2}E.
\]
\end{corollary}

\begin{proof}
Combine Lemma~\ref{lem:rho-to-cubic} with Lemma~\ref{lem:acceleration}.
\end{proof}

\section{Exact level-set calculus along a geodesic}

Let \(x(t)\), \(t\in[0,1]\), be a unit-time Hessian geodesic with constant kinetic energy
\[
    E:=\|\dot x(t)\|_{x(t)}^2.
\]
Fix \(y\), and set
\[
    v(t):=V(x(t),y),
    \qquad
    v_0:=v(0).
\]

\begin{lemma}[Second variation identity]\label{lem:second-variation}
Along a Hessian geodesic,
\[
    v''(t)=2E-\ip{\ddot x(t)}{U(x(t),y)}_{x(t)}.
\]
Consequently, if
\[
    v'(0)\le -(E+v_0),
\]
then for every \(t\in[0,1]\),
\[
\begin{aligned}
    v(t)
    &\le
    v_0-tv_0-Et(1-t) \\
    &\quad+
    \int_0^t(t-s)
    \left|
    \ip{\ddot x(s)}{U(x(s),y)}_{x(s)}
    \right|ds.
\end{aligned}
\]
In particular,
\[
    v(1)
    \le
    \int_0^1(1-t)
    \left|
    \ip{\ddot x(t)}{U(x(t),y)}_{x(t)}
    \right|dt.
\]
\end{lemma}

\begin{proof}
Differentiating
\[
    v(t)=\ip{x(t)-y}{g(x(t))-g(y)}
\]
gives
\[
    v'(t)=
    \ip{\dot x(t)}{g(x(t))-g(y)}+
    \ip{x(t)-y}{H(x(t))\dot x(t)}.
\]
Differentiating again yields
\[
\begin{aligned}
    v''(t)
    &=
    \ip{\ddot x(t)}{g(x(t))-g(y)}
    +2\dot x(t)^\top H(x(t))\dot x(t) \\
    &\quad+
    \ip{\ddot x(t)}{H(x(t))(x(t)-y)}
    +D^3\phi(x(t))[\dot x(t),\dot x(t),x(t)-y].
\end{aligned}
\]
Since \(E=\dot x(t)^\top H(x(t))\dot x(t)\), the middle term is \(2E\). The Hessian geodesic equation is equivalent to
\[
    \ip{\ddot x(t)}{w}_{x(t)}
    =
    -\frac12D^3\phi(x(t))[\dot x(t),\dot x(t),w]
    \qquad \text{for every }w.
\]
Taking \(w=x(t)-y\), we get
\[
    D^3\phi(x(t))[\dot x(t),\dot x(t),x(t)-y]
    =
    -2\ip{\ddot x(t)}{x(t)-y}_{x(t)}.
\]
Therefore
\[
\begin{aligned}
    v''(t)
    &=
    2E+
    \ip{\ddot x(t)}{g(x(t))-g(y)}
    +\ip{\ddot x(t)}{x(t)-y}_{x(t)}
    -2\ip{\ddot x(t)}{x(t)-y}_{x(t)} \\
    &=
    2E-
    \ip{\ddot x(t)}{(x(t)-y)-H(x(t))^{-1}(g(x(t))-g(y))}_{x(t)} \\
    &=
    2E-
    \ip{\ddot x(t)}{U(x(t),y)}_{x(t)}.
\end{aligned}
\]
This proves the second variation identity.

For the integral estimate, use the exact formula
\[
    v(t)=v_0+tv'(0)+\int_0^t(t-s)v''(s)\,ds.
\]
Substituting the second variation identity gives
\[
    v(t)=v_0+tv'(0)+Et^2-
    \int_0^t(t-s)\ip{\ddot x(s)}{U(x(s),y)}_{x(s)}ds.
\]
Using \(v'(0)\le -(E+v_0)\), we obtain
\[
\begin{aligned}
    v(t)
    &\le
    v_0-tv_0-Et(1-t)
    +
    \int_0^t(t-s)
    \left|
    \ip{\ddot x(s)}{U(x(s),y)}_{x(s)}
    \right|ds.
\end{aligned}
\]
Setting \(t=1\) proves the endpoint estimate.
\end{proof}


\begin{lemma}[Abstract sublevel bootstrap, signed-work form]
Let \(x(t)\), \(t\in[0,1]\), be a unit-time Hessian geodesic with constant
kinetic energy
\[
    E:=\|\dot x(t)\|_{x(t)}^2.
\]
Fix \(y\), and define
\[
    v(t):=V(x(t),y),
    \qquad
    v_0:=v(0).
\]
Assume \(v_0>0\) and
\[
    v'(0)\le -(E+v_0).
\]
Let
\[
    W(t):=\langle \ddot x(t),U(x(t),y)\rangle_{x(t)}.
\]
Suppose there is a nondecreasing function \(F:[0,\infty)\to[0,\infty)\)
such that, whenever \(v(s)\le v_0\) on \([0,t]\),
\[
    \int_0^t (t-s)(-W(s))_+\,ds
    \le
    tF(v_0).
\]
If
\[
    F(v_0)<v_0,
\]
then
\[
    v(t)<v_0
    \qquad
    \forall t\in(0,1],
\]
and
\[
    v(1)\le F(v_0).
\]
\end{lemma}

\begin{proof}
For any \(t\in[0,1]\),
\[
    v(t)
    =
    v_0+t v'(0)+\int_0^t(t-s)v''(s)\,ds.
\]
Using
\[
    v''(s)=2E-W(s),
\]
we get
\[
    v(t)
    =
    v_0+t v'(0)+Et^2-\int_0^t(t-s)W(s)\,ds.
\]
The descent condition gives
\[
    t v'(0)\le -t(E+v_0),
\]
so
\[
    v(t)
    \le
    v_0-tv_0-Et(1-t)
    -
    \int_0^t(t-s)W(s)\,ds.
\]
Since
\[
    -W(s)\le (-W(s))_+,
\]
we have
\[
    -\int_0^t(t-s)W(s)\,ds
    \le
    \int_0^t(t-s)(-W(s))_+\,ds.
\]
Thus, on any interval where \(v(s)\le v_0\),
\[
    v(t)
    \le
    v_0-tv_0-Et(1-t)+tF(v_0).
\]
Dropping the nonpositive term \(-Et(1-t)\),
\[
    v(t)
    \le
    v_0-t\bigl(v_0-F(v_0)\bigr).
\]
If \(F(v_0)<v_0\), this is strictly less than \(v_0\) for all \(t>0\)
as long as the estimate is valid. A first-exit argument then shows that
the path cannot leave the initial sublevel set. Taking \(t=1\) gives
\[
    v(1)\le F(v_0).
\]
\end{proof}



\begin{lemma}[Abstract sublevel bootstrap]\label{lem:bootstrap}
Let \(x(t)\), \(t\in[0,1]\), be a unit-time Hessian geodesic with constant
kinetic energy
\[
    E:=\|\dot x(t)\|_{x(t)}^2.
\]
Fix \(y\), and define
\[
    v(t):=V(x(t),y),
    \qquad
    v_0:=v(0).
\]
Assume \(v_0>0\) and the descent condition
\[
    v'(0)\le -(E+v_0).
\]
Suppose there is a nondecreasing function \(F:[0,\infty)\to[0,\infty)\)
such that, for all \(s\in[0,1]\),
\[
    \frac12
    \left|
    \langle \ddot x(s),U(x(s),y)\rangle_{x(s)}
    \right|
    \le
    F(v(s)).
\]
If
\[
    F(v_0)<v_0,
\]
then
\[
    v(t)<v_0
    \qquad
    \forall t\in(0,1],
\]
and
\[
    v(1)\le F(v_0).
\]
\end{lemma}

\begin{proof}
For any \(t\in[0,1]\), the exact integral identity gives
\[
    v(t)
    =
    v_0+t v'(0)+\int_0^t(t-s)v''(s)\,ds.
\]
Using
\[
    v''(s)
    =
    2E-\langle \ddot x(s),U(x(s),y)\rangle_{x(s)},
\]
we obtain
\[
\begin{aligned}
    v(t)
    &=
    v_0+t v'(0)+Et^2
    -
    \int_0^t(t-s)
    \langle \ddot x(s),U(x(s),y)\rangle_{x(s)}\,ds.
\end{aligned}
\]
The descent condition implies
\[
    t v'(0)\le -t(E+v_0).
\]
Therefore
\[
\begin{aligned}
    v(t)
    &\le
    v_0-tv_0-Et(1-t)
    +
    \int_0^t(t-s)
    \left|
    \langle \ddot x(s),U(x(s),y)\rangle_{x(s)}
    \right|ds.
\end{aligned}
\]

Now consider any interval \([0,t]\) on which
\[
    v(s)\le v_0
    \qquad
    \forall s\in[0,t].
\]
Since \(F\) is nondecreasing, the pointwise hypothesis gives
\[
    \left|
    \langle \ddot x(s),U(x(s),y)\rangle_{x(s)}
    \right|
    \le
    2F(v(s))
    \le
    2F(v_0)
    \qquad
    \forall s\in[0,t].
\]
Hence
\[
\begin{aligned}
    v(t)
    &\le
    v_0-tv_0-Et(1-t)
    +
    2F(v_0)\int_0^t(t-s)\,ds  \\
    &=
    v_0-tv_0-Et(1-t)
    +
    F(v_0)t^2.
\end{aligned}
\]
Dropping the nonpositive term \(-Et(1-t)\) and using \(t^2\le t\), we get
\[
    v(t)
    \le
    v_0-tv_0+tF(v_0)
    =
    v_0-t\bigl(v_0-F(v_0)\bigr).
\]

If \(F(v_0)<v_0\), this is strictly less than \(v_0\) for every \(t>0\),
as long as the estimate is valid. We now use a first-exit argument. Since
\[
    v'(0)\le -(E+v_0)<0,
\]
the path initially enters the strict sublevel set. Suppose, for contradiction,
that the path exits the initial sublevel set. Let \(\tau>0\) be the first time
such that
\[
    v(\tau)=v_0,
    \qquad
    v(t)<v_0
    \quad
    \text{for }0<t<\tau.
\]
Then the preceding estimate is valid on \([0,\tau]\), and gives
\[
    v(\tau)
    \le
    v_0-\tau\bigl(v_0-F(v_0)\bigr)
    <
    v_0,
\]
contradicting \(v(\tau)=v_0\). Therefore no such first exit time exists, and
\[
    v(t)<v_0
    \qquad
    \forall t\in(0,1].
\]

Taking \(t=1\) in the estimate gives
\[
    v(1)
    \le
    v_0-\bigl(v_0-F(v_0)\bigr)
    =
    F(v_0).
\]
\end{proof}



\section{Self-concordant specialization}

We use the convention for standard self-concordance
\[
    |D^3\phi(x)[a,b,b]|
    \le
    2\|a\|_x\|b\|_x^2.
\]
The next lemma records the standard implication from self-concordance to
exponential Hessian stability. This makes the use of \(\rho(r)=e^{2r}\) in this
section self-contained.

\begin{lemma}[Self-concordance implies exponential Hessian stability]\label{lem:sc-to-stability}
Assume the more general cubic bound
\[
    |D^3\phi(x)[a,b,b]|
    \le
    L\|a\|_x\|b\|_x^2
\]
for all \(x,a,b\). Then the Hessian metric is exponentially stable with modulus
\[
    \rho(r)=e^{Lr}.
\]
That is, for all \(x,y\in\Omega\),
\[
    e^{-Ld_R(x,y)}H(x)
    \preceq
    H(y)
    \preceq
    e^{Ld_R(x,y)}H(x).
\]
In particular, standard self-concordance corresponds to \(L=2\), hence
\(\rho(r)=e^{2r}\).
\end{lemma}

\begin{proof}
Let \(R:=d_R(x,y)\), and let \(\gamma:[0,R]\to\Omega\) be a unit-speed
geodesic from \(x\) to \(y\). Fix an arbitrary vector \(u\), and define
\[
    Q(s):=u^\top H(\gamma(s))u=\|u\|_{\gamma(s)}^2.
\]
Then
\[
    Q'(s)=D^3\phi(\gamma(s))[\gamma'(s),u,u].
\]
By the cubic bound and the unit-speed parametrization,
\[
    |Q'(s)|
    \le
    L\|\gamma'(s)\|_{\gamma(s)}\|u\|_{\gamma(s)}^2
    =LQ(s).
\]
Thus
\[
    -L\le \frac{d}{ds}\log Q(s)\le L
\]
whenever \(Q(s)>0\). Integrating from \(0\) to \(R\) gives
\[
    e^{-LR}Q(0)\le Q(R)\le e^{LR}Q(0).
\]
Equivalently,
\[
    e^{-Ld_R(x,y)}u^\top H(x)u
    \le
    u^\top H(y)u
    \le
    e^{Ld_R(x,y)}u^\top H(x)u.
\]
Since this holds for every \(u\), the Loewner-order statement follows.
\end{proof}

By Lemma~\ref{lem:sc-to-stability}, standard self-concordance gives
\(\rho(r)=e^{2r}\). Lemma~\ref{lem:acceleration} also gives
\[
    \|\ddot x(t)\|_{x(t)}\le E.
\]

\begin{theorem}[Self-concordant level-set recursion]\label{thm:sc-recursion}
Assume standard self-concordance. Let \(x(t)\) be a unit-time Hessian geodesic with energy \(E\), and define \(v(t)=V(x(t),y)\), \(v_0=v(0)\). Suppose
\[
    v'(0)\le -(E+v_0),
    \qquad
    E\le d(x(0),y).
\]
Define
\[
    A_{\rm sc}(v):=2(\cosh\sqrt v-1),
\]
\[
    D_{\rm sc}(v):=2v+A_{\rm sc}(v)^2,
\]
and
\[
    F_{\rm sc}(v):=\frac12D_{\rm sc}(v)A_{\rm sc}(v).
\]
If
\[
    F_{\rm sc}(v_0)<v_0,
\]
then
\[
    v(t)<v_0\qquad\forall t\in(0,1],
\]
and
\[
    v(1)\le F_{\rm sc}(v_0).
\]
Moreover,
\[
    F_{\rm sc}(v)=v^2+O(v^3)
    \qquad (v\downarrow0).
\]
\end{theorem}

\begin{proof}
Since \(\rho(r)=e^{2r}\), Lemma~\ref{lem:U-bound} gives
\[
    \|U(z,y)\|_z
    \le
    2(\cosh d_R(z,y)-1).
\]
Using Lemma~\ref{lem:distance-V}, \(d_R(z,y)^2\le V(z,y)\), we get
\[
    \|U(z,y)\|_z
    \le
    A_{\rm sc}(V(z,y)).
\]
At the initial point, Lemma~\ref{lem:parallelogram} and \(E\le d(x(0),y)\) give
\[
    E
    \le
    2v_0+\|U(x(0),y)\|_{x(0)}^2
    \le
    D_{\rm sc}(v_0).
\]
By Lemma~\ref{lem:acceleration}, \(\|\ddot x(t)\|_{x(t)}\le E\). Hence Lemma~\ref{lem:bootstrap} applies.
This proves the sublevel invariance and the one-step bound.

Finally,
\[
    A_{\rm sc}(v)=2(\cosh\sqrt v-1)=v+\frac{v^2}{12}+O(v^3),
\]
so
\[
    F_{\rm sc}(v)
    =
    \frac12\bigl(2v+A_{\rm sc}(v)^2\bigr)A_{\rm sc}(v)
    =v^2+O(v^3).
\]
\end{proof}

\begin{corollary}[Quadratic convergence on a certified level set]
Under the assumptions of Theorem~\ref{thm:sc-recursion}, the recurrence
\[
    v_{k+1}\le F_{\rm sc}(v_k)
\]
is quadratically convergent whenever the iterates remain in a level set where \(F_{\rm sc}(v)<v\). More precisely, if \(0<\delta<0.683412\ldots\), then \(F_{\rm sc}(v)<v\) for all \(0<v\le\delta\), and
\[
    v_{k+1}\le C_\delta v_k^2,
    \qquad
    C_\delta:=\sup_{0<v\le\delta}\frac{F_{\rm sc}(v)}{v^2}<\infty.
\]
In particular, the simple choice \(\delta=1/4\) is valid, and
\[
    C_{1/4}\approx 1.155.
\]
\end{corollary}


\section{General modulus version}

The preceding argument only uses three inputs: a bound on \(U\), the identity controlling \(d\) by \(V\) and \(U\), and an acceleration bound. The following theorem records the resulting general-modulus statement.

\begin{theorem}[General-modulus level-set recursion]\label{thm:general-modulus}
Assume \(\rho\)-Hessian stability. Define
\[
    A_\rho(v):=J_\rho(\sqrt v).
\]
Assume that Hessian geodesics satisfy the acceleration estimate
\[
    \|\ddot x(t)\|_{x(t)}\le \kappa_\rho E
\]
for some constant \(\kappa_\rho\).  In particular, if \(\rho\) is compatible, then
Corollary~\ref{cor:rho-acceleration} allows the choice \(\kappa_\rho=L_\rho/2\).
Let \(x(t)\) be a unit-time Hessian geodesic with energy \(E\), and define \(v(t)=V(x(t),y)\), \(v_0=v(0)\). Suppose
\[
    v'(0)\le -(E+v_0),
    \qquad
    E\le d(x(0),y).
\]
Define
\[
    F_\rho(v):=\frac{\kappa_\rho}{2}\bigl(2v+A_\rho(v)^2\bigr)A_\rho(v).
\]
If
\[
    F_\rho(v_0)<v_0,
\]
then the geodesic step remains in the initial sublevel set,
\[
    v(t)<v_0 \qquad \forall t\in(0,1],
\]
and
\[
    v(1)\le F_\rho(v_0).
\]
\end{theorem}

\begin{proof}
By Lemma~\ref{lem:U-bound} and Lemma~\ref{lem:distance-V},
\[
    \|U(x,y)\|_x\le A_\rho(V(x,y)).
\]
By Lemma~\ref{lem:parallelogram}, if \(E\le d(x,y)\), then
\[
    E\le 2V(x,y)+A_\rho(V(x,y))^2.
\]
Thus the hypotheses of Lemma~\ref{lem:bootstrap} holds. 
The conclusion follows directly from Lemma~\ref{lem:bootstrap}.
\end{proof}

\begin{corollary}[Local quadratic rate from a differentiable modulus]
Suppose, in addition, that
\[
    \rho(r)=1+L_\rho r+O(r^2)
    \qquad (r\downarrow0),
\]
and that \(\kappa_\rho=L_\rho/2\). Then
\[
    F_\rho(v)=\frac{L_\rho^2}{4}v^2+O(v^{5/2}).
\]
\end{corollary}

\begin{proof}
Since
\[
    K_\rho(r)-K_\rho(r)^{-1}=L_\rho r+O(r^2),
\]
we have
\[
    A_\rho(v)=J_\rho(\sqrt v)=\frac{L_\rho}{2}v+O(v^{3/2}).
\]
Substituting this into the definition of \(F_\rho\) gives
\[
    F_\rho(v)
    =\frac{L_\rho/2}{2}\left(2v+O(v^2)\right)
      \left(\frac{L_\rho}{2}v+O(v^{3/2})\right)
    =\frac{L_\rho^2}{4}v^2+O(v^{5/2}).
\]
\end{proof}

\begin{corollary}[Exponential modulus]
For \(\rho(r)=e^{Lr}\), one has
\[
    A_\rho(v)=\frac{4}{L}\left(\cosh\left(\frac{L}{2}\sqrt v\right)-1\right).
\]
If \(\kappa_\rho=L/2\), then
\[
    F_L(v)=
    \frac{L}{4}
    \left(2v+
    \frac{16}{L^2}\left(\cosh\left(\frac{L}{2}\sqrt v\right)-1\right)^2
    \right)
    \frac{4}{L}\left(\cosh\left(\frac{L}{2}\sqrt v\right)-1\right).
\]
For \(L=2\), this recovers \(F_{\rm sc}\).
\end{corollary}

\end{document}
